\documentclass[10pt,a4paper]{amsart}
\usepackage{amsmath,amssymb,amsthm,graphicx}
% Retained minimal preamble for the Stacks Project chapter flat.tex (More on Flatness).
% Verbatim excerpt of the chapter preamble: Xy-pic setup (native preamble lines 14-18) and the native
% theorem declarations with the mathematical macros (native preamble lines 154-217). Copied mechanically;
% nothing added, removed or redefined. The wrapper class and the amsmath/amssymb/amsthm/graphicx base
% packages are supplied by the assembler, so the class/externaldocument part of the native preamble is omitted.
% For commutative diagrams we use Xy-pic
\usepackage[all]{xy}

% We use 2cell for 2-commutative diagrams.
\xyoption{2cell}
% Theorem environments.
%
\theoremstyle{plain}
\newtheorem{theorem}[subsection]{Theorem}
\newtheorem{proposition}[subsection]{Proposition}
\newtheorem{lemma}[subsection]{Lemma}

\theoremstyle{definition}
\newtheorem{definition}[subsection]{Definition}
\newtheorem{example}[subsection]{Example}
\newtheorem{exercise}[subsection]{Exercise}
\newtheorem{situation}[subsection]{Situation}

\theoremstyle{remark}
\newtheorem{remark}[subsection]{Remark}
\newtheorem{remarks}[subsection]{Remarks}

\numberwithin{equation}{subsection}

% Macros
%
\def\lim{\mathop{\mathrm{lim}}\nolimits}
\def\colim{\mathop{\mathrm{colim}}\nolimits}
\def\Spec{\mathop{\mathrm{Spec}}}
\def\Hom{\mathop{\mathrm{Hom}}\nolimits}
\def\Ext{\mathop{\mathrm{Ext}}\nolimits}
\def\SheafHom{\mathop{\mathcal{H}\!\mathit{om}}\nolimits}
\def\SheafExt{\mathop{\mathcal{E}\!\mathit{xt}}\nolimits}
\def\Sch{\mathit{Sch}}
\def\Mor{\mathop{\mathrm{Mor}}\nolimits}
\def\Ob{\mathop{\mathrm{Ob}}\nolimits}
\def\Sh{\mathop{\mathit{Sh}}\nolimits}
\def\NL{\mathop{N\!L}\nolimits}
\def\CH{\mathop{\mathrm{CH}}\nolimits}
\def\proetale{{pro\text{-}\acute{e}tale}}
\def\etale{{\acute{e}tale}}
\def\QCoh{\mathit{QCoh}}
\def\Ker{\mathop{\mathrm{Ker}}}
\def\Im{\mathop{\mathrm{Im}}}
\def\Coker{\mathop{\mathrm{Coker}}}
\def\Coim{\mathop{\mathrm{Coim}}}

% Boxtimes
%
\DeclareMathSymbol{\boxtimes}{\mathbin}{AMSa}{"02}

%
% Macros for moduli stacks/spaces
%
\def\QCohstack{\mathcal{QC}\!\mathit{oh}}
\def\Cohstack{\mathcal{C}\!\mathit{oh}}
\def\Spacesstack{\mathcal{S}\!\mathit{paces}}
\def\Quotfunctor{\mathrm{Quot}}
\def\Hilbfunctor{\mathrm{Hilb}}
\def\Curvesstack{\mathcal{C}\!\mathit{urves}}
\def\Polarizedstack{\mathcal{P}\!\mathit{olarized}}
\def\Complexesstack{\mathcal{C}\!\mathit{omplexes}}
% \Pic is the operator that assigns to X its picard group, usage \Pic(X)
% \Picardstack_{X/B} denotes the Picard stack of X over B
% \Picardfunctor_{X/B} denotes the Picard functor of X over B
\def\Pic{\mathop{\mathrm{Pic}}\nolimits}
\def\Picardstack{\mathcal{P}\!\mathit{ic}}
\def\Picardfunctor{\mathrm{Pic}}
\def\Deformationcategory{\mathcal{D}\!\mathit{ef}}

% --- appended 2026-10-12 (Chief): the project own definitions of the reference/slogan/history
% display environments, copied mechanically from the shared Stacks preamble lines 8-12
% (\usepackage{verbatim} plus the three \newenvironment lines, which the chapter preamble keeps
% verbatim from the project preamble). Nothing added, removed or redefined: these environments
% carry no mathematical content and the project itself typesets them as comments. Without these
% three lines a counted statement whose own text uses them fails with Environment slogan undefined.
\usepackage{verbatim}
\newenvironment{reference}{\comment}{\endcomment}
\newenvironment{slogan}{\comment}{\endcomment}
\newenvironment{history}{\comment}{\endcomment}
\usepackage{cleveref}
\crefname{originalreference}{original reference}{original references}
\Crefname{originalreference}{Original reference}{Original references}
\makeatletter
\newlabel{situation-existence}{{0CTC}{source}{}{sourcepointer-1}{}}
\newlabel{situation-existence@cref}{{[originalreference][0][]0CTC}{[1][0][]source}{}{}{}}
\makeatother
\crefname{originallemma}{original lemma}{original lemmas}
\Crefname{originallemma}{Original lemma}{Original lemmas}
\makeatletter
\newlabel{lemma-compute-what-it-should-be}{{0CTD}{source}{}{sourcepointer-2}{}}
\newlabel{lemma-compute-what-it-should-be@cref}{{[originallemma][0][]0CTD}{[1][0][]source}{}{}{}}
\makeatother
\crefname{originallemma}{original lemma}{original lemmas}
\Crefname{originallemma}{Original lemma}{Original lemmas}
\makeatletter
\newlabel{lemma-relative-pseudo-coherence}{{0CTF}{source}{}{sourcepointer-3}{}}
\newlabel{lemma-relative-pseudo-coherence@cref}{{[originallemma][0][]0CTF}{[1][0][]source}{}{}{}}
\makeatother
\crefname{originallemma}{original lemma}{original lemmas}
\Crefname{originallemma}{Original lemma}{Original lemmas}
\makeatletter
\newlabel{lemma-compute-over-affine}{{0CTG}{source}{}{sourcepointer-4}{}}
\newlabel{lemma-compute-over-affine@cref}{{[originallemma][0][]0CTG}{[1][0][]source}{}{}{}}
\makeatother
\title{On a neighbourhood of the closed subset where the derived limit lives, the complex has a single cohomology sheaf of finite presentation reducing correctly along the thickenings}
\author{The Stacks Project Authors}
\date{stacks-project:flat.tex:local-snapshot-20261011; GNU FDL 1.2 or later}
\begin{document}
\maketitle
\noindent\textit{Editorial note.} All delivered mathematical text is the author's own, copied from the pinned Stacks Project chapter file flat.tex; no statement or proof step was rewritten, reordered or removed.\par\medskip
\noindent\textit{Editorial note.} The preamble is the chapter's own retained minimal preamble (sources/stacks-flat/retained-preamble.tex), a verbatim excerpt of the native Xy-pic setup and of the native theorem declarations and mathematical macros, followed by the project own definitions of the reference, slogan and history display environments (verbatim package plus three newenvironment lines) that the retained file carries at its end and that the project itself typesets as comments. No macro was invented or redefined, and the wrapper class and the amsmath/amssymb/amsthm/graphicx base packages are supplied by the assembler.\par\medskip
\noindent\textit{Editorial note.} The counted statement and proof contain no cross reference, no citation and no included file, so the delivered body is self contained.\par\medskip
\noindent\textit{Editorial note.} Canonical location of the counted proposition: Stacks tag 0CTH (sources/stacks-flat/pinned-original-source.tex line 8181).\par\medskip
\noindent\textit{Allowed standard prerequisites (editorial).}\begin{itemize}
\item The existence situation and the complexes K and F\_n. Same chapter, flat.tex, Situation labelled situation-existence (source line 8036) and Lemma labelled lemma-compute-what-it-should-be (source line 8062), both invoked in the statement. Fix the objects K, F\_n and the canonical map K -> F\_n which gives the meaning of part (3).
\item Relative pseudo-coherence and the affine computation of the cohomology of K. Same chapter, Lemma labelled lemma-relative-pseudo-coherence (source line 8128) invoked at proof lines 8209-8214, and Lemma labelled lemma-compute-over-affine (source line 8145) invoked at proof lines 8216-8225. Supply the bounded above complex of finite free P-modules representing Ri\_*K and the computation of its cohomology after tensoring with A\_n.
\end{itemize}\medskip
\ifdefined\hypertarget\hypertarget{sourcepointer-1}{}\fi
\noindent\textit{Source reference.} Original reference 0CTC: sources/stacks-flat/pinned-original-source.tex, source lines 8036--8036. Same-chapter pointer into the author environment situation at source line 8036; invoked at proof line None of the counted proof. It is retained only as a navigation pointer; every load-bearing step is inside the delivered span.\par\medskip
\ifdefined\hypertarget\hypertarget{sourcepointer-2}{}\fi
\noindent\textit{Source reference.} Original lemma 0CTD: sources/stacks-flat/pinned-original-source.tex, source lines 8062--8062. Same-chapter pointer into the author environment lemma at source line 8062; invoked at proof line None of the counted proof. It is retained only as a navigation pointer; every load-bearing step is inside the delivered span.\par\medskip
\ifdefined\hypertarget\hypertarget{sourcepointer-3}{}\fi
\noindent\textit{Source reference.} Original lemma 0CTF: sources/stacks-flat/pinned-original-source.tex, source lines 8128--8128. Same-chapter pointer into the author environment lemma at source line 8128; invoked at proof line 8209 of the counted proof. It is retained only as a navigation pointer; every load-bearing step is inside the delivered span.\par\medskip
\ifdefined\hypertarget\hypertarget{sourcepointer-4}{}\fi
\noindent\textit{Source reference.} Original lemma 0CTG: sources/stacks-flat/pinned-original-source.tex, source lines 8145--8145. Same-chapter pointer into the author environment lemma at source line 8145; invoked at proof line 8216 of the counted proof. It is retained only as a navigation pointer; every load-bearing step is inside the delivered span.\par\medskip
\begin{lemma}
\label{lemma-finitely-presented}
In Situation \ref{situation-existence} let $K$ be as in
Lemma \ref{lemma-compute-what-it-should-be}.
Denote $X_0 \subset X$ the closed subset
consisting of points lying over the closed subset
$\Spec(A_1) = \Spec(A_2) = \ldots$ of $\Spec(A)$.
There exists an open $W \subset X$ containing $X_0$
such that
\begin{enumerate}
\item $H^i(K)|_W$ is zero unless $i = 0$,
\item $\mathcal{F} = H^0(K)|_W$ is of finite presentation, and
\item $\mathcal{F}_n = \mathcal{F} \otimes_{\mathcal{O}_X} \mathcal{O}_{X_n}$.
\end{enumerate}
\end{lemma}

\begin{proof}
Fix $n \geq 1$. By construction there is a canonical map
$K \to \mathcal{F}_n$ in $D_\QCoh(\mathcal{O}_X)$
and hence a canonical map $H^0(K) \to \mathcal{F}_n$
of quasi-coherent sheaves. This explains the meaning of part (3).

\medskip\noindent
Let $x \in X_0$ be a point.
We will find an open neighbourhood $W$
of $x$ such that (1), (2), and (3) are true. Since $X_0$ is quasi-compact
this will prove the lemma. Let $U \subset X$ be an affine open
neighbourhood of $x$. Say $U = \Spec(B)$.
Choose a surjection $P \to B$ with $P$ smooth over $A$.
By Lemma \ref{lemma-relative-pseudo-coherence}
and the definition of relative pseudo-coherence
there exists a bounded above complex $F^\bullet$
of finite free $P$-modules representing
$Ri_*K$ where $i : U \to \Spec(P)$ is the closed
immersion induced by the presentation.
Let $M_n$ be the $B$-module corresponding to $\mathcal{F}_n|_U$.
By Lemma \ref{lemma-compute-over-affine}
$$
H^i(F^\bullet \otimes_A A_n) =
\left\{
\begin{matrix}
0 & \text{if} & i \not = 0 \\
M_n & \text{if} & i = 0
\end{matrix}
\right.
$$
Let $i$ be the maximal index such that $F^i$ is nonzero.
If $i \leq 0$, then (1), (2), and (3) are true.
If not, then $i > 0$ and we see that the rank of the map
$$
F^{i - 1} \to F^i
$$
in the point $x$ is maximal. Hence in an open neighbourhood
of $x$ inside $\Spec(P)$ the rank is maximal. Thus after replacing
$P$ by a principal localization we may assume that the displayed
map is surjective. Since $F^i$ is finite free we may choose
a splitting $F^{i - 1} = F' \oplus F^i$. Then we may
replace $F^\bullet$ by the complex
$$
\ldots \to F^{i - 2} \to F' \to 0 \to \ldots
$$
and we win by induction on $i$.
\end{proof}

\end{document}
